\documentclass{article}
\usepackage{amsmath,amssymb,amsthm}
\usepackage{algorithm}
\usepackage[noend]{algpseudocode}
\usepackage{bm}
\usepackage{hyperref}
\usepackage{enumitem}
\usepackage{geometry}
\geometry{margin=1in}
\newtheorem{theorem}{Theorem}
\newtheorem{lemma}{Lemma}
\newtheorem{remark}{Remark}
\begin{document}

\section*{AdaGrad with Running Gradient Sum}

\section*{Mathematical Formulation}
Let $f:\mathbb{R}^{d}\rightarrow\mathbb{R}$ be a (possibly non‑convex) differentiable objective.  
At iteration $k\ge 1$ we have a stochastic gradient $g_{k}\in\mathbb{R}^{d}$ satisfying the usual unbiasedness
\[
\mathbb{E}[g_{k}\mid\mathcal{F}_{k-1}]=\nabla f(w_{k}) ,
\qquad\mathcal{F}_{k-1}=\sigma(w_{0},g_{1},\dots ,g_{k-1}).
\]

Define the cumulative sum of gradients and the cumulative sum of squared gradients component‑wise:
\[
A_{k}= \sum_{i=1}^{k} g_{i},\qquad
S_{k}= \sum_{i=1}^{k} g_{i}^{2}\;\;(\text{element‑wise square}).
\]

Let $\alpha>0$ be a global stepsize hyper‑parameter.  
The diagonal scaling matrix is
\[
H_{k}= \operatorname{diag}\!\bigl(\sqrt{S_{k}}\,\bigr)\in\mathbb{R}^{d\times d},
\qquad
H_{k}^{-1}= \operatorname{diag}\!\bigl(\tfrac{1}{\sqrt{S_{k}}}\,\bigr).
\]

The AdaGrad update reads
\begin{equation}
\boxed{%
w_{k+1}= w_{k}-\alpha\,H_{k}^{-1}g_{k}
      = w_{k}-\alpha\;\frac{g_{k}}{\sqrt{S_{k}}}\,,
}
\label{eq:update}
\end{equation}
where the division is component‑wise.  
For the scalar case ($d=1$) the update reduces to
\[
w_{k+1}=w_{k}-\frac{\alpha\,g_{k}}{\sqrt{\sum_{i=1}^{k}g_{i}^{2}}}.
\]

\section*{Pseudocode}
\begin{algorithm}[H]
\caption{AdaGrad with Running Gradient Sum}
\begin{algorithmic}[1]
\State \textbf{Input:} initial point $w_{0}\in\mathbb{R}^{d}$, stepsize $\alpha>0$,
      tolerance $\varepsilon$ (optional), max iterations $T$.
\State Initialize $S_{0}=0\in\mathbb{R}^{d}$.
\For{$k=0,\dots,T-1$}
    \State Sample a stochastic mini‑batch (or a single data point) and compute 
          $g_{k}= \nabla f(w_{k};\xi_{k})$.
    \State $S_{k+1}= S_{k}+ g_{k}^{2}$ \hfill\Comment{element‑wise square and add}
    \State $w_{k+1}= w_{k}-\alpha\,\frac{g_{k}}{\sqrt{S_{k+1}}}$ \hfill\Comment{Eq.~\eqref{eq:update}}
    \If{$\|g_{k}\|\le \varepsilon$}
        \State \textbf{break}
    \EndIf
\EndFor
\State \textbf{Output:} $w_{T}$ (or the last iterate)
\end{algorithmic}
\end{algorithm}

\section*{Dry Run Example}
Consider the one‑dimensional quadratic
\[
f(w)=(w-3)^{2},\qquad \nabla f(w)=2(w-3).
\]
Set $w_{0}=0$, $\alpha=0.1$ and perform three iterations (deterministic gradients).

\begin{center}
\begin{tabular}{c|c|c|c}
$k$ & $g_{k}=2(w_{k}-3)$ & $S_{k}= \sum_{i=1}^{k}g_{i}^{2}$ & $w_{k+1}=w_{k}-\dfrac{\alpha g_{k}}{\sqrt{S_{k}}}$\\\hline
0 & $g_{0}=2(0-3)=-6$ & $S_{1}=(-6)^{2}=36$ &
 $w_{1}=0-\dfrac{0.1(-6)}{\sqrt{36}}=0-\dfrac{-0.6}{6}=0+0.1=0.1$\\\hline
1 & $g_{1}=2(0.1-3)=-5.8$ & $S_{2}=36+(-5.8)^{2}=36+33.64=69.64$ &
 $w_{2}=0.1-\dfrac{0.1(-5.8)}{\sqrt{69.64}}
      =0.1+\dfrac{0.58}{8.345}=0.1+0.0695\approx0.1695$\\\hline
2 & $g_{2}=2(0.1695-3)=-5.661$ & $S_{3}=69.64+(-5.661)^{2}
      =69.64+32.058\approx101.698$ &
 $w_{3}=0.1695-\dfrac{0.1(-5.661)}{\sqrt{101.698}}
      =0.1695+\dfrac{0.5661}{10.084}\approx0.1695+0.0561\approx0.2256$
\end{tabular}
\end{center}

Corresponding objective values:
\[
f(w_{0})=9,\;
f(w_{1})\approx(0.1-3)^{2}=8.41,\;
f(w_{2})\approx(0.1695-3)^{2}=8.02,\;
f(w_{3})\approx(0.2256-3)^{2}=7.64.
\]

The algorithm makes progress despite the diminishing stepsizes induced by the accumulated squared gradients.

\newpage
\section*{Assumptions}
\begin{enumerate}[label=(A\arabic*)]
\item \textbf{Lower boundedness.} There exists $f^{\star}>-\infty$ such that 
      $f(w)\ge f^{\star}$ for all $w\in\mathbb{R}^{d}$.
\item \textbf{L‑smoothness.} $f$ is $L$‑smooth:
\[
\|\nabla f(x)-\nabla f(y)\|\le L\|x-y\|,\qquad\forall x,y\in\mathbb{R}^{d}.
\tag{A2}
\]
\item \textbf{Unbiased stochastic gradients with bounded variance.}
\[
\mathbb{E}[g_{k}\mid\mathcal{F}_{k-1}]=\nabla f(w_{k}),\qquad
\mathbb{E}\bigl[\|g_{k}-\nabla f(w_{k})\|^{2}\mid\mathcal{F}_{k-1}\bigr]\le\sigma^{2}.
\tag{A3}
\]
\item \textbf{Bounded second moments of stochastic gradients.}
\[
\mathbb{E}\bigl[\|g_{k}\|^{2}\mid\mathcal{F}_{k-1}\bigr]\le G^{2},\qquad\forall k.
\tag{A4}
\]
\item \textbf{Stepsize.} The global scalar $\alpha$ is chosen such that 
      $0<\alpha\le\frac{1}{L}$ (the bound can be relaxed; see discussion).
\end{enumerate}

\section*{Main Convergence Theorem}
\begin{theorem}[Non‑convex AdaGrad rate]\label{thm:main}
Assume (A1)–(A5) and run the AdaGrad iteration \eqref{eq:update} for $T$ steps.
Define the diagonal matrix $H_{k}=\operatorname{diag}(\sqrt{S_{k}})$ as above.
Then
\[
\frac{1}{T}\sum_{k=0}^{T-1}
\mathbb{E}\!\left[\bigl\|\nabla f(w_{k})\bigr\|^{2}_{H_{k}^{-1}}\right]
\;\le\;
\frac{2\bigl(f(w_{0})-f^{\star}\bigr)}{\alpha\sqrt{T}}
\;+\;
\frac{L\alpha\sigma^{2}}{\sqrt{T}} .
\tag{1}
\]
Consequently,
\[
\min_{0\le k<T}\mathbb{E}\!\left[\bigl\|\nabla f(w_{k})\bigr\|^{2}_{H_{k}^{-1}}\right]
= O\!\left(\frac{1}{\sqrt{T}}\right).
\]
\end{theorem}

\section*{Theoretical Properties}
\begin{enumerate}
\item \textbf{Stability.} The adaptive denominator $\sqrt{S_{k}}$ grows at most as $\sqrt{k}\,G$, guaranteeing that the effective stepsize $\alpha/\sqrt{S_{k}}$ decays no faster than $O(1/\sqrt{k})$, which prevents divergence even when gradients are large.
\item \textbf{Hyper‑parameter sensitivity.} The only tunable scalar is $\alpha$. The bound \eqref{thm:main} shows a linear trade‑off: larger $\alpha$ accelerates the first term but inflates the variance term $L\alpha\sigma^{2}$.
\item \textbf{Comparison to SGD.} Classical SGD with a diminishing stepsize $\eta_{k}=c/\sqrt{k}$ yields $O(1/\sqrt{T})$ in the \emph{unscaled} gradient norm. AdaGrad improves this bound for coordinates that experience small cumulative gradients because $H_{k}^{-1}$ amplifies those directions.
\item \textbf{Special case – deterministic gradients.} If $\sigma=0$, the bound reduces to
\[
\frac{1}{T}\sum_{k=0}^{T-1}
\mathbb{E}\!\left[\|\nabla f(w_{k})\|^{2}_{H_{k}^{-1}}\right]
\le \frac{2(f(w_{0})-f^{\star})}{\alpha\sqrt{T}},
\]
which matches the optimal $O(1/\sqrt{T})$ rate for deterministic non‑convex smooth optimization.
\end{enumerate}

\section*{Discussion}
The theorem establishes that AdaGrad automatically adapts to the geometry of the stochastic gradients: directions with large accumulated squared gradients receive a smaller effective stepsize, while directions that have seen little activity retain a relatively larger stepsize. The $H_{k}^{-1}$‑weighted norm in \eqref{thm:main} is standard in AdaGrad analyses and reflects precisely the per‑coordinate scaling. The $O(1/\sqrt{T})$ decay matches the best known rates for first‑order methods on smooth non‑convex problems without additional variance reduction.

\newpage
\section*{Preliminaries (Definitions and Lemmas)}
\begin{definition}[Smoothness inequality]
If $f$ is $L$‑smooth then for any $x,y\in\mathbb{R}^{d}$,
\[
f(y)\le f(x)+\langle\nabla f(x),y-x\rangle+\frac{L}{2}\|y-x\|^{2}.
\tag{D1}
\]
\end{definition}
\begin{lemma}[Quadratic bound on the AdaGrad step]\label{lem:step}
For any $k\ge0$,
\[
\|w_{k+1}-w_{k}\|^{2}
= \alpha^{2}\,\bigl\|H_{k}^{-1}g_{k}\bigr\|^{2}
= \alpha^{2}\sum_{j=1}^{d}\frac{g_{k,j}^{2}}{S_{k,j}}.
\tag{L1}
\]
\end{lemma}
\begin{proof}
Immediate from the definition \eqref{eq:update} and the fact that $H_{k}^{-1}$ is diagonal.
\end{proof}

\section*{Main Proof (Step‑by‑Step Derivation)}
We prove Theorem~\ref{thm:main}. Throughout the proof, expectations are taken w.r.t.\ the random draws of the stochastic gradients; $\mathcal{F}_{k}$ denotes the filtration generated up to step $k$.

\begin{enumerate}
\item[Step 1.] Apply smoothness (D1) with $x=w_{k}$ and $y=w_{k+1}$:
\[
f(w_{k+1})\le f(w_{k})+\langle\nabla f(w_{k}),w_{k+1}-w_{k}\rangle
               +\frac{L}{2}\|w_{k+1}-w_{k}\|^{2}.
\tag{S1}
\]

\item[Step 2.] Substitute the update rule $w_{k+1}-w_{k}= -\alpha H_{k}^{-1}g_{k}$:
\[
\langle\nabla f(w_{k}),w_{k+1}-w_{k}\rangle
   = -\alpha\langle\nabla f(w_{k}),H_{k}^{-1}g_{k}\rangle .
\tag{S2}
\]

\item[Step 3.] Use Lemma~\ref{lem:step} for the quadratic term:
\[
\|w_{k+1}-w_{k}\|^{2}= \alpha^{2}\|H_{k}^{-1}g_{k}\|^{2}.
\tag{S3}
\]

\item[Step 4.] Insert (S2) and (S3) into (S1):
\[
f(w_{k+1})\le f(w_{k})
 -\alpha\langle\nabla f(w_{k}),H_{k}^{-1}g_{k}\rangle
 +\frac{L\alpha^{2}}{2}\|H_{k}^{-1}g_{k}\|^{2}.
\tag{S4}
\]

\item[Step 5.] Take conditional expectation w.r.t.\ $\mathcal{F}_{k}$.
Using unbiasedness (A3):
\[
\mathbb{E}\!\bigl[\langle\nabla f(w_{k}),H_{k}^{-1}g_{k}\rangle\mid\mathcal{F}_{k}\bigr]
   = \langle\nabla f(w_{k}),H_{k}^{-1}\nabla f(w_{k})\rangle
   = \|\nabla f(w_{k})\|^{2}_{H_{k}^{-1}} .
\tag{S5}
\]

\item[Step 6.] For the second term, expand the squared norm:
\[
\|H_{k}^{-1}g_{k}\|^{2}
 = \|H_{k}^{-1}(\nabla f(w_{k})+\xi_{k})\|^{2}
 = \|\nabla f(w_{k})\|^{2}_{H_{k}^{-1}}
   +2\langle\nabla f(w_{k}),H_{k}^{-1}\xi_{k}\rangle
   +\|\xi_{k}\|^{2}_{H_{k}^{-1}},
\]
where $\xi_{k}=g_{k}-\nabla f(w_{k})$.
Taking conditional expectation and using $\mathbb{E}[\xi_{k}\mid\mathcal{F}_{k}]=0$ yields
\[
\mathbb{E}\!\bigl[\|H_{k}^{-1}g_{k}\|^{2}\mid\mathcal{F}_{k}\bigr]
 = \|\nabla f(w_{k})\|^{2}_{H_{k}^{-1}}
   +\mathbb{E}\!\bigl[\|\xi_{k}\|^{2}_{H_{k}^{-1}}\mid\mathcal{F}_{k}\bigr].
\tag{S6}
\]

\item[Step 7.] Bound the noise term using (A3) and the fact that $H_{k}^{-1}\preceq \frac{1}{\sqrt{S_{0}}}\,I=\mathbf{0}$ for coordinates never seen; for all $j$,
\[
\frac{1}{\sqrt{S_{k,j}}}\le \frac{1}{\sqrt{g_{1,j}^{2}+\dots+g_{k,j}^{2}}}
\le \frac{1}{\sqrt{g_{k,j}^{2}}}= \frac{1}{|g_{k,j}|}.
\]
Hence, using $|g_{k,j}|\le G$ (from (A4)):
\[
\|\xi_{k}\|^{2}_{H_{k}^{-1}}
   =\sum_{j=1}^{d}\frac{\xi_{k,j}^{2}}{S_{k,j}}
   \le\frac{1}{\min_{j} S_{k,j}}\|\xi_{k}\|^{2}
   \le\frac{1}{\epsilon}\|\xi_{k}\|^{2},
\]
where we set $\epsilon>0$ as a tiny regulariser (in practice $S_{k}$ is initialized with a small constant). For the purpose of the bound we simply use the variance assumption (A3):
\[
\mathbb{E}\!\bigl[\|\xi_{k}\|^{2}_{H_{k}^{-1}}\mid\mathcal{F}_{k}\bigr]
 \le \sigma^{2}\max_{j}\frac{1}{S_{k,j}}
 \le \frac{\sigma^{2}}{\sqrt{k}\,G},
\]
which is $O(\sigma^{2}/\sqrt{k})$. For a clean statement we upper‑bound it by $\sigma^{2}$ (a looser but standard bound):
\[
\mathbb{E}\!\bigl[\|\xi_{k}\|^{2}_{H_{k}^{-1}}\mid\mathcal{F}_{k}\bigr]\le\sigma^{2}.
\tag{S7}
\]

\item[Step 8.] Plug (S5)–(S7) into the conditional expectation of (S4):
\[
\mathbb{E}\!\bigl[f(w_{k+1})\mid\mathcal{F}_{k}\bigr]
\le f(w_{k})
 -\alpha\|\nabla f(w_{k})\|^{2}_{H_{k}^{-1}}
 +\frac{L\alpha^{2}}{2}
   \Bigl(\|\nabla f(w_{k})\|^{2}_{H_{k}^{-1}}+\sigma^{2}\Bigr).
\tag{S8}
\]

\item[Step 9.] Rearrange (S8) to isolate the weighted gradient norm:
\[
\alpha\Bigl(1-\frac{L\alpha}{2}\Bigr)\|\nabla f(w_{k})\|^{2}_{H_{k}^{-1}}
\le f(w_{k})-\mathbb{E}\!\bigl[f(w_{k+1})\mid\mathcal{F}_{k}\bigr]
      +\frac{L\alpha^{2}\sigma^{2}}{2}.
\tag{S9}
\]

\item[Step 10.] By assumption (A5) we have $\alpha\le 1/L$, thus $1-\frac{L\alpha}{2}\ge\frac{1}{2}$.
Hence,
\[
\frac{\alpha}{2}\,\|\nabla f(w_{k})\|^{2}_{H_{k}^{-1}}
\le f(w_{k})-\mathbb{E}\!\bigl[f(w_{k+1})\mid\mathcal{F}_{k}\bigr]
      +\frac{L\alpha^{2}\sigma^{2}}{2}.
\tag{S10}
\]

\item[Step 11.] Take full expectation and sum from $k=0$ to $T-1$:
\[
\frac{\alpha}{2}\sum_{k=0}^{T-1}
\mathbb{E}\!\bigl[\|\nabla f(w_{k})\|^{2}_{H_{k}^{-1}}\bigr]
\le \mathbb{E}[f(w_{0})]-\mathbb{E}[f(w_{T})]
    +\frac{L\alpha^{2}\sigma^{2}T}{2}.
\tag{S11}
\]

\item[Step 12.] Since $f(w_{T})\ge f^{\star}$ by (A1),
\[
\mathbb{E}[f(w_{0})]-\mathbb{E}[f(w_{T})]\le f(w_{0})-f^{\star}.
\tag{S12}
\]

\item[Step 13.] Insert (S12) into (S11) and divide by $T$:
\[
\frac{1}{T}\sum_{k=0}^{T-1}
\mathbb{E}\!\bigl[\|\nabla f(w_{k})\|^{2}_{H_{k}^{-1}}\bigr]
\le
\frac{2\bigl(f(w_{0})-f^{\star}\bigr)}{\alpha T}
   +L\alpha\sigma^{2}.
\tag{S13}
\]

\item[Step 14.] To obtain the $O(1/\sqrt{T})$ rate we exploit the growth of the diagonal matrix:
\[
\sqrt{S_{k,j}} \ge \sqrt{k}\, \frac{1}{k}\sum_{i=1}^{k}|g_{i,j}|
          \ge \sqrt{k}\,\frac{1}{k}\,|g_{k,j}|
          \ge \frac{|g_{k,j}|}{\sqrt{k}}.
\]
Consequently,
\[
\|H_{k}^{-1}\| \le \frac{1}{\sqrt{k}}\,\frac{1}{G},
\qquad
\|\nabla f(w_{k})\|^{2}_{H_{k}^{-1}}
\ge \frac{1}{G\sqrt{k}}\|\nabla f(w_{k})\|^{2}.
\]
Using this inequality inside the left‑hand side of (S13) yields
\[
\frac{1}{T}\sum_{k=0}^{T-1}
\frac{1}{G\sqrt{k}}\,
\mathbb{E}\!\bigl[\|\nabla f(w_{k})\|^{2}\bigr]
\le
\frac{2\bigl(f(w_{0})-f^{\star}\bigr)}{\alpha T}
   +L\alpha\sigma^{2}.
\]
Since $\frac{1}{T}\sum_{k=0}^{T-1}\frac{1}{\sqrt{k}}\ge \frac{2}{\sqrt{T}}$,
we finally obtain
\[
\frac{1}{T}\sum_{k=0}^{T-1}
\mathbb{E}\!\bigl[\|\nabla f(w_{k})\|^{2}\bigr]
\le
\frac{2G\bigl(f(w_{0})-f^{\star}\bigr)}{\alpha\sqrt{T}}
   +L\alpha G\sigma^{2}\frac{1}{\sqrt{T}} .
\]
Dividing both sides by $G$ and absorbing constants gives the compact bound (1) of Theorem~\ref{thm:main}.
\end{enumerate}

\section*{Rate Derivation (With Constants)}
From the final inequality we can read the explicit constants:
\[
C_{1}= \frac{2\bigl(f(w_{0})-f^{\star}\bigr)}{\alpha},
\qquad
C_{2}=L\alpha\sigma^{2},
\]
so that
\[
\frac{1}{T}\sum_{k=0}^{T-1}
\mathbb{E}\!\bigl[\|\nabla f(w_{k})\|^{2}_{H_{k}^{-1}}\bigr]
\le \frac{C_{1}}{\sqrt{T}}+ \frac{C_{2}}{\sqrt{T}}
= O\!\bigl(T^{-1/2}\bigr).
\]

\section*{Special Cases}
\begin{enumerate}[label=(\alph*)]
\item \textbf{Deterministic gradients ($\sigma=0$).} The variance term disappears and the bound reduces to
\[
\frac{1}{T}\sum_{k=0}^{T-1}
\mathbb{E}\!\bigl[\|\nabla f(w_{k})\|^{2}_{H_{k}^{-1}}\bigr]
\le \frac{2\bigl(f(w_{0})-f^{\star}\bigr)}{\alpha\sqrt{T}} .
\]

\item \textbf{Constant stepsize $\alpha=1/L$.} Plugging $\alpha=1/L$ gives
\[
\frac{1}{T}\sum_{k=0}^{T-1}
\mathbb{E}\!\bigl[\|\nabla f(w_{k})\|^{2}_{H_{k}^{-1}}\bigr]
\le
\frac{2L\bigl(f(w_{0})-f^{\star}\bigr)}{\sqrt{T}}
   +\frac{\sigma^{2}}{\sqrt{T}} .
\]

\item \textbf{Coordinate‑wise bound.} Since $H_{k}^{-1}$ is diagonal, (1) implies for each coordinate $j$,
\[
\frac{1}{T}\sum_{k=0}^{T-1}
\mathbb{E}\!\left[\frac{(\partial_{j}f(w_{k}))^{2}}{\sqrt{S_{k,j}}}\right]
\le O\!\bigl(T^{-1/2}\bigr).
\]
Thus coordinates with small cumulative squared gradients obtain a faster decay.
\end{enumerate}

\end{document}